% chapters/14/contoh.tex
\section{Dua Contoh Terapan}
\label{sec:contoh-bab-14}

\subsection{Contoh Pertama: Menghitung Ukuran Gambar pada Scene}

\begin{contohbox}
Sebuah gambar pemain berukuran 128 kali 128 piksel diimpor dengan
\istilah{Pixels Per Unit} bernilai 100. Tinggi pemain pada scene adalah

\[
128 \div 100 = 1{,}28 \text{ satuan}.
\]

Apabila tinggi pemain yang dikehendaki adalah 2 satuan, maka nilai
\istilah{Pixels Per Unit} yang diperlukan adalah

\[
128 \div 2 = 64 .
\]
\end{contohbox}

Perhitungan itu menjelaskan mengapa \istilah{Pixels Per Unit} perlu diatur
sebelum levelnya disusun. Ukuran gambar pada layar tidak ditentukan oleh
berkasnya, melainkan oleh perbandingan antara ukuran berkas dan nilai
pengaturan itu.

Perhatikan pula bahwa seluruh aset sebaiknya memakai nilai yang sama. Apabila
gambar pemain bernilai 64 sedangkan gambar rintangan bernilai 100, maka ukuran
keduanya tidak dapat dibandingkan secara langsung, dan Anda akan menghabiskan
waktu untuk menyesuaikan ukurannya satu per satu.

\subsection{Contoh Kedua: Menelusuri Animasi yang Tidak Berganti}

\begin{contohbox}
Seorang mahasiswa melaporkan bahwa animasi pemainnya tidak pernah berganti
dari diam menjadi berlari, walaupun pemainnya sudah bergerak. Program berjalan
tanpa pesan kesalahan apa pun.
\end{contohbox}

Penelusurannya dilakukan menurut urutan berikut.

\begin{langkah}
  \item \textbf{Periksa jendela \istilah{Animator}.} Apakah kedua keadaannya
        ada, dan apakah transisinya tergambar? Ternyata keduanya ada, sehingga
        masalahnya bukan pada pengaturan animasinya.
  \item \textbf{Periksa nilainya pada jendela \istilah{Animator} ketika scene
        berjalan.} Jendela itu menampilkan nilai setiap parameter pada saat
        itu. Ternyata nilai kecepatannya tetap nol, walaupun pemainnya
        bergerak.
  \item \textbf{Periksa kode yang menetapkan nilainya.} Ternyata nama parameter
        pada kode, yaitu \code{speed}, berbeda dari nama pada
        \istilah{Animator}, yaitu \code{Speed}. Perbedaan besar-kecil huruf
        itu tidak diperiksa, dan nilainya ditetapkan pada parameter yang tidak
        ada.
  \item \textbf{Perbaiki, lalu jalankan kembali.} Animasinya berganti
        sebagaimana mestinya.
\end{langkah}

Tiga hal dapat dipetik dari penelusuran itu. Pertama, jendela
\istilah{Animator} dapat menampilkan nilai parameter ketika scene berjalan,
dan itulah alat pemeriksa yang paling berguna di sini. Kedua, kesalahan
penamaan tidak menimbulkan pesan apa pun, sehingga hanya dapat ditemukan
dengan memeriksa nilainya. Ketiga, memakai peubah tetap untuk nama parameter,
sebagaimana pada skrip Daftar~\ref{lst:ch14-animasi}, mencegah kesalahan itu
sejak awal.
